\input{../tex-inputs/latex-leseansicht-vorspann}

\section[Gustav Schwarzkopf an Arthur Schnitzler, 21. 7. 1899]{L04300 Gustav Schwarzkopf an Arthur Schnitzler, 21. 7. 1899}
\nopagebreak\mylabel{L04300v}
\rehead{ }\normalsize\beginnumbering\briefempfaengerindex{Schnitzler, Arthur@\textsc{Schnitzler, Arthur}!zzzSchwarzkopf, Gustav@\emph{von Gustav Schwarzkopf}!1899-07-212@{21. 7. 1899}|(be}
\toendnotes[C]{\smallbreak\pagebreak[2]}
\correspDesc{Versand  durch Gustav Schwarzkopf am 21. 7. 1899 in Wien
\newline{}Übermittlung  am 21. 7. 1899 in Wien
\newline{}Erhalt  durch Arthur Schnitzler im Zeitraum [22. 7. 1899 – 26. 7. 1899?] in Velden am Wörthersee}\toendnotes[C]{\smallbreak}
\Standort{CUL, Schnitzler, B 96a.}
\physDesc{Postkarte, 549 Zeichen
\newline{}Handschrift: schwarze Tinte, deutsche Kurrent
\newline{}Versand: 1) Stempel: »\nobreak{}\oindex{I., Innere Stadt@\textbf{I., Innere Stadt}, \emph{Verwaltungsgebiet}|pwk}Wien 1/1 1, 21. 7. 99, 8–9N\nobreak{}«.   2) Stempel: »\nobreak{}\oindex{Velden am Wörthersee@\textbf{Velden am Wörthersee}|pwk}Velden am Wörtersee, \textcolor{gray}{2}\textcolor{gray}{×}\textcolor{gray}{. 7.} 99, 8.F\nobreak{}«. }\toendnotes[C]{\smallbreak}\pstart{}{\pb}Herrn \textsc{D\textsuperscript{r} Arthur Schnitzler}\pend{}\pstart{}\textsc{Velden} am Wörther-See\oindex{Velden am Wörthersee@\textbf{Velden am Wörthersee}|pw}\pend{}\pstart{}\textsc{Pension Pundschu\oindex{Pension Pundschu@\textbf{Pension Pundschu}, \emph{Hotel}|pw}}\pend{}{\bigskip}\vspace{1em}
\pstart
           \raggedleft{}{\pb}21. 7. 99\pend
           \vspace{0.5em}
\pstart
           Lieber Arthur, ich will Ihnen nur den Empfang Ihrer \label{K_L04300-1v}\edtext{Karte}{\lemma{\textnormal{\emph{Karte}}}\Cendnote{\textnormal{{XXXX ref} XXXX}}}\label{K_L04300-1}
               beſtätigen. Zu erzählen weiß ich nichts. Ich hab kein i{\geminationn}eres Erlebnis gehabt, es geht auch{ }ſonſt nichts vor. Nicht einmal die fällige »Fackel\pwindex{Fackel@\emph{Die Fackel}|pw}« ist erſchienen. Es iſt Ihr wahrſcheinlich
               zu heiß. — Mit \uline{wem} haben Sie eine{ }ſchöne Radpartie
               gemacht?\pend
           
\pstart
           Ich bitte Sie mich Ihrer Frau Mama\pwindex{Schnitzler, Louise 8.\,7.\,1840 Kőszeg – 9.\,9.\,1911 Wien@\textsc{Schnitzler, Louise} (8.\,7.\,1840 Kőszeg – 9.\,9.\,1911 Wien)|pwv} und Ihrer Frau Schweſter\pwindex{Hajek, Gisela 20.\,12.\,1867 Wien – 3.\,2.\,1953 Cambridge@\textsc{Hajek, Gisela} (20.\,12.\,1867 Wien – 3.\,2.\,1953 Cambridge)|pwv} beſtens zu empfehlen und B.
                  H.\pwindex{Beer-Hofmann, Richard 11.\,7.\,1866 Wien – 26.\,9.\,1945 New York City@\textsc{Beer-Hofmann, Richard} (11.\,7.\,1866 Wien – 26.\,9.\,1945 New York City), \emph{Schriftsteller}|pw} zu grüßen, we{\geminationn} Sie ihn nächſtens doch{ }ſehen{ }ſollten, wie ich annehmen will.\pend
           
\pstart
           Herzlichſt{\\[\baselineskip]}Ihr{\\[\baselineskip]}\spacefill\mbox{Gustav Sch.}\pend
           \leftskip=0em{}\selectlanguage{ngerman}\endnumbering\briefempfaengerindex{Schnitzler, Arthur@\textsc{Schnitzler, Arthur}!zzzSchwarzkopf, Gustav@\emph{von Gustav Schwarzkopf}!1899-07-212@{21. 7. 1899}|)be}\mylabel{L04300h}
\begin{anhang}
\end{anhang}\newcommand{\dateiname}{L04300}\newcommand{\titel}{Gustav Schwarzkopf an Arthur Schnitzler, 21. 7. 1899}\newcommand{\editorInnen}{Herausgegeben von Jahnke, SelmaMüller, Martin Anton}\input{../tex-inputs/latex-leseansicht-abspann}
